\clearpage
\chapter*{Fe de erratas de la exposición n\textsuperscript{o} 149}
\addcontentsline{toc}{chapter}{Fe de erratas de la exposición n\textsuperscript{o} 149}
\setcounter{footnote}{0}
\src{298}

\noindent\textbf{1. Teorema de dualidad para los haces algebraicos
coherentes.}

\begin{center}
[Seminario Bourbaki, t. 9, 1956/57, n\textsuperscript{o} 149, 25 p.]
\end{center}

\noindent\textbf{Observación sobre la página 14.} --- Como señalé en mi
conferencia en el Congreso Internacional de Matemáticos de 1958,\footnote{
Grothendieck (Alexander). --- The cohomology theory of abstract algebraic
varieties, \emph{Proceedings of the International Congress of Mathematicians}
[1958, Edinburgh], p. 103--118. --- Cambridge, at the University Press, 1960.}
las cuestiones planteadas aquí están ahora completamente resueltas.

El lector encontrará más información sobre la dualidad de los haces
coherentes en \emph{loco citato}\footnotemark[1], p. 112--115, en EGA III,
2\textsuperscript{a} parte (en preparación) y, mientras tanto, en SGA 1962.
Un tratamiento más sistemático se encontrará en un capítulo posterior de EGA
(capítulo IX en el plan previsto).
